\section{Non-Gaussian margins}
\label{sec:stress}

\Cref{tab:stress} separates two kinds of statement.
Alignment~(\cref{thm:alignment}) and reliability~(\cref{thm:reliability}) are second-moment identities.
They hold for every law with finite second moments; Monte Carlo agreement only checks that sample moments converge.
The informative column is the Gaussian $(z,r^2)$ map for log-loss, Brier score, AUC and mutual information.

\begin{table}[t]
  \centering
  \caption{Stress battery.
           Second-moment identities hold by construction.
           The $(z,r^2)$ column is a Monte Carlo check of the Gaussian map at the same second moments.}
  \label{tab:stress}
  \begin{tabular}{@{}lcc@{}}
    \toprule
    Law & Gaussian $(z,r^2)$ map & Second-moment identities \\
    \midrule
    Gaussian            & matches & by construction \\
    Repeated team       & matches & by construction$^\ast$ \\
    Student $t_3$       & fails   & by construction \\
    Student $t_5$       & fails   & by construction \\
    Gamma margins       & fails   & by construction \\
    Bivariate Poisson   & fails   & by construction \\
    Ordinal cut         & fails   & by construction \\
    Five per cent contamination & fails & by construction \\
    \bottomrule
  \end{tabular}\\[0.5em]
  {\footnotesize $^\ast$On the misaligned repeated-team design the \emph{sample} reliability is $0.733$ against population theory $0.753$ ($z=-3.2$).
   That discrepancy is estimation under dependent matches, not a failure of~\cref{thm:reliability}.}
\end{table}

The $(z,r^2)$ reduction matches for the Gaussian law and for the repeated-team law, which is marginally Gaussian.
It fails for $t_3$, $t_5$, gamma margins, bivariate Poisson and the ordinal cut.
Contamination is mixed across metrics and does not restore the Gaussian map.

Three predictive questions under non-Gaussian margins are left open: whether logistic $\beta_A+\beta_B=0$ still tracks best-linear-predictor alignment; whether $D$ loses nothing to the pair in out-of-sample log-loss when $\beta\propto c$; and whether the $1/(2n)$ cost becomes the Takeuchi correction under misspecification.
They are not claimed as results of this draft.

The repeated-team simulation recovers a latent $\eta_S/\eta_N$ by Gaussian maximum likelihood with a Woodbury likelihood, using Wald intervals at nominal level $95\%$.
Coverage is $0.983$ at a $16$-team design ($240$ matches, $60$ replicates) and $0.967$ at a $24$-team design ($1{,}104$ matches, $30$ replicates).
The population reliability identity assumes independent signal and noise and i.i.d.\ draws.
Team effects induce dependence across matches that share a club, so a sample reliability computed as if matches were independent can deviate from the population value even when~\cref{thm:reliability} holds.
Match-level applications are reported in the companion empirical paper~\parencite{brownPEFempirical}.
